\section{The symmetric power and Hilbert--Chow morphism}
\source{fga-explained:page-0169}\source{fga-explained:page-0170}\source{fga-explained:page-0171}\source{fga-explained:page-0172}\source{fga-explained:page-0173}\source{fga-explained:page-0174}\source{fga-explained:page-0175}\source{fga-explained:page-0176}\source{fga-explained:page-0177}\source{fga-explained:page-0178}\source{fga-explained:page-0179}\source{fga-explained:page-0180}\source{fga-explained:page-0181}\source{fga-explained:page-0182}\source{fga-explained:page-0183}\source{fga-explained:page-0184}\source{fga-explained:page-0185}\source{fga-explained:page-0186}\source{fga-explained:page-0187}\source{fga-explained:page-0188}

Fix an algebraically closed field $k$ of characteristic zero and a
quasi-projective $k$-scheme $X$.  The component $\Hilb^n(X)$ parametrizes
zero-dimensional subschemes of length $n$.

\begin{definition}[symmetric power]\label{fga-def-7-1-1}\dcref{7.1.1}\source{fga-explained:page-0169}\uses{}
$X^{(n)}=X^n/S_n$ is the symmetric power; its points are effective zero-cycles
$\sum n_i[p_i]$ of degree $n$.
\end{definition}\begin{theorem}[Hilbert--Chow]\label{fga-thm-7-1-2}\dcref{7.1.2}\source{fga-explained:page-0169}\uses{fga-def-7-1-1,fga-def-5-1}
There is a morphism $\rho:\Hilb^n(X)\to X^{(n)}$ sending a finite flat family
to the cycle whose coefficient at $p$ is the length of the local fiber at $p$.
It is an isomorphism for nonsingular curves and is birational for nonsingular
surfaces.
\end{theorem}
\begin{proof}
The determinant of multiplication on the finite locally free algebra of a
family gives the elementary symmetric functions, hence a map to the quotient
$X^n/S_n$.  On a smooth curve, ideals of colength $n$ are uniquely determined
by their divisor, giving an inverse.  On a smooth surface the locus of distinct
points is dense and the map is an isomorphism there, so it is birational.
\end{proof}
\begin{lemma}\label{fga-lem-7-1-4}\dcref{7.1.4}\source{fga-explained:page-0170}\uses{fga-thm-7-1-2}
The punctual Hilbert scheme of length $n$ supported at a smooth point of a
surface is connected and has the expected dimension; its open locus of reduced
clusters is dense in the global Hilbert scheme.
\end{lemma}
\begin{lemma}\label{fga-lem-7-1-7}\dcref{7.1.7}\source{fga-explained:page-0171}\uses{fga-thm-6-1-19}
For a finite subscheme $Z\subset X$, the Zariski tangent space to $\Hilb^n(X)$
at $[Z]$ is $\operatorname{Hom}(\mathcal I_Z,\OO_Z)$.
\end{lemma}
\begin{lemma}\label{fga-lem-7-1-8}\dcref{7.1.8}\source{fga-explained:page-0171}\uses{fga-lem-7-1-7}
If $X$ is a smooth surface, the tangent dimension at every cluster is at least
$2n$ and equals $2n$ on the reduced locus.
\end{lemma}
\begin{theorem}\label{fga-thm-7-1-13}\dcref{7.1.13}\source{fga-explained:page-0173}\uses{fga-lem-7-1-7,fga-lem-7-1-8}
For a nonsingular surface $X$, $\Hilb^n(X)$ is nonsingular of dimension $2n$
and irreducible.
\end{theorem}
\begin{proof}
Induct on $n$ using the Hilbert--Chow fibers and the punctual calculation.
The tangent-space bound and the dimension of the dense reduced locus force
smoothness; connected punctual fibers and the birational map to the irreducible
symmetric power give irreducibility.
\end{proof}

\section{Examples and stratifications}
\begin{lemma}\label{fga-lem-7-1-11}\dcref{7.1.11}\source{fga-explained:page-0172}\uses{fga-thm-7-1-2}
The Hilbert scheme of points on a smooth curve is the symmetric power, and for
$X=\mathbb A^2$ the fiber over $n[0]$ is the punctual Hilbert scheme.
\end{lemma}
\begin{theorem}\label{fga-thm-7-1-14}\dcref{7.1.14}\source{fga-explained:page-0174}\uses{fga-thm-7-1-13,fga-thm-7-1-2}
The Hilbert--Chow morphism for a smooth surface is a projective resolution of
the singularities of $X^{(n)}$.
\end{theorem}
\begin{lemma}\label{fga-lem-7-2-1}\dcref{7.2.1}\source{fga-explained:page-0175}\uses{fga-def-7-1-1}
The strata of $X^{(n)}$ indexed by partitions of $n$ have codimension determined
by the number of distinct support points; their inverse images stratify
$\Hilb^n(X)$.
\end{lemma}
\begin{theorem}\label{fga-thm-7-2-3}\dcref{7.2.3}\source{fga-explained:page-0175}\uses{fga-lem-7-2-1,fga-thm-7-1-2}
For a smooth surface, the open stratum of distinct points is dense and its
complement has codimension at least two; the Hilbert--Chow map is an isomorphism
over that open stratum.
\end{theorem}
\begin{definition}\label{fga-def-7-4-1}\dcref{7.4.1}\source{fga-explained:page-0177}\uses{fga-def-5-1}
The nested Hilbert scheme and the support stratification record a flag
$Z_1\subset Z_2\subset\cdots\subset Z_n$ of clusters and the successive
support points; these are locally closed functors.
\end{definition}
\begin{proposition}\label{fga-prop-7-4-4}\dcref{7.4.4}\source{fga-explained:page-0178}\uses{fga-def-7-4-1,fga-thm-7-1-2}
The support strata are locally closed, their closures are ordered by refinement
of partitions, and the Hilbert--Chow morphism respects this stratification.
\end{proposition}

\section{Betti numbers and the Heisenberg algebra}
\begin{definition}\label{fga-def-7-5-1}\dcref{7.5.1}\source{fga-explained:page-0180}\uses{}
For a smooth projective surface $X$, write $P_t(Y)=\sum b_i(Y)t^i$ for the
Poincare polynomial.  The generating series of the Betti numbers of
$\Hilb^n(X)$ is formed over all $n\ge0$.
\end{definition}\begin{theorem}\label{fga-thm-7-5-5}\dcref{7.5.5}\source{fga-explained:page-0182}\uses{fga-def-7-5-1,fga-thm-7-1-13}
Gottsche's formula is
\[
\sum_{n\ge0}P_t(\Hilb^n(X))q^n=
\prod_{m\ge1}\frac{\prod_{i\ \mathrm{odd}}(1+t^iq^m)^{b_i(X)}}
{\prod_{i\ \mathrm{even}}(1-t^iq^m)^{b_i(X)}}.
\]
\end{theorem}
\begin{definition}\label{fga-def-7-6-1}\dcref{7.6.1}\source{fga-explained:page-0185}\uses{fga-def-7-5-1}
The Nakajima correspondences define creation and annihilation operators on
$\bigoplus_n H^*(\Hilb^n(X))$; their commutator is the intersection pairing on
$H^*(X)$.
\end{definition}
\begin{theorem}\label{fga-thm-7-6-2}\dcref{7.6.2}\source{fga-explained:page-0187}\uses{fga-def-7-6-1,fga-thm-7-5-5}
These operators furnish a representation of the Heisenberg algebra, and the
Fock-space character of the representation is the product in
\cref{fga-thm-7-5-5}.
\end{theorem}
